\section{Critical masses and the analytic inputs}\label{sec:inputs}
Our geometric arguments use positive sums of walk and polygon weights.
This section fixes their endpoint conventions and states the precise
analytic estimates proved in the companion papers. A free endpoint always
means a sum over the actual possible endpoint ports from one fixed source;
it never introduces a second translation sum.

\subsection{Ports, strips and local endpoint changes}
A port is the midpoint of an edge of the triangular tiling. A port path
follows honeycomb edges, with an initial and a final half-edge, and visits
each dual vertex at most once. A nonempty port path has distinct endpoint
ports; the empty height-zero bridge below is the sole exception. Its
length counts visited vertices and its
weight is \(w(\gamma)=\rho^{|\gamma|}\), with
\(\rho=(2+\sqrt2)^{-1/2}\). The trace of a port path includes its two
ports. Removing these ports changes its diameter by a bounded constant.
Unrestricted vertex paths keep the edge-count convention of the introduction.

A horizontal band is the region between two successive horizontal lines
of tiling edges. Its Euclidean height is \(d_0=\sqrt3/2\).
A strict bridge of height \(h\) goes from a bottom port to a top port of
the open \(h\)-band strip and has no other visit to the boundary.
An arch has two distinct ports on the same boundary. Because ports carry
no weight and adjacent open strips have disjoint vertices, concatenating
strict bridges multiplies their weights exactly. Every horizontal seam
port is a translate of the initial port by a symmetry of the triangular
tiling, with its inward half-edge pointing into the next band. Thus each
translated continuation has exactly the same admissible shapes and weights;
there is no residual port type or parity state in this concatenation. Tangential displacements
of height-\(h\) bridges lie in \(h/2+\Z\). Write
\[
 b_h(x)=\sum_{\gamma:0\to(x,h)}w(\gamma),\qquad
 B_h=\sum_xb_h(x),\qquad B_0=1.
\]
We write \(\Prob_h^{\rm br}\) for the probability measure obtained by
dividing these bridge weights by \(B_h\). Let \(K_h(j)\) be the
arch mass between bottom ports separated by the nonzero integer \(j\),
and put \(A_h=\sum_{j\ne0}K_h(j)\) and
\(m_h=\sum_{j>0}jK_h(j)\).

A weak bridge is a vertex path whose endpoints attain its minimum and
maximum heights; intermediate visits to these levels are permitted.
The following bounded modification will be used whenever a geometric
cut produces a vertex endpoint.
\begin{lemma}[Endpoint extension]\label{lem:port-extension}
A weak bridge in any lattice normal direction can be extended to a strict
port bridge between the first lattice boundary lines outside its height
range. The height and length change by \(O(1)\), the weight ratio is
bounded above and below by positive constants, and the inverse operation
has bounded multiplicity after recording its endpoint types.
\end{lemma}
\begin{proof}
At a lower endpoint whose triangle center is nearest its lower side,
append the vertical half-edge to that side's port. At the other triangle
type, first append its fixed downward edge and then the vertical half-edge.
The reflected operation applies at the upper endpoint. Every newly added
vertex lies strictly outside the old height range, so it was unvisited.
The two additions are disjoint for positive macroscopic span. Bounded
spans have only finitely many endpoint types and use the same construction;
a zero-span piece is trivial. At most three extra vertex factors occur
relative to the original edge count. Removing the recorded one- or
two-edge endpoint patterns recovers the original path, proving all bounds.
Rotations and reflections give the other lattice normal directions.
\end{proof}

Throughout, a bound \(F(h)\le h^{a+o(1)}\) means that for each
\(\epsilon>0\), \(F(h)\le C_\epsilon h^{a+\epsilon}\) for all
sufficiently large \(h\). The corresponding lower bound has the analogous
meaning. Constants may depend on declared fixed geometry and on exponent
slack. Uniform statements retain one choice of constant over their full
stated parameter range.

\subsection{Strip flux and nesting}
Set \(c=\cos(3\pi/8)\). The first input is the positive boundary
consequence of the parafermionic observable, together with its quantitative
strip refinement. We use the complete statements of
\cite[Theorem 1.1 and Lemma 2.1]{StripCompanion}.
\begin{theorem}[Boundary flux and strip mass]\label{thm:01-strip}\label{lem:01-boundary}
For every finite simply connected union \(D\) of triangular tiles with
simple polygonal boundary and each boundary source port \(a\),
\[
 \sum_{b\in\partial D}\sum_{\gamma:a\to b}
       \rho^{|\gamma|}e^{3iW(\gamma)/8}=1,
\]
where \(W\) is total signed turning from the initial inward direction to
the final outward direction. If \(D\) is convex, its unsigned exit mass
is at most \(1/c\). For each \(h\ge1\), all strip sums defined above
are finite and
\[
 cA_h+B_h=1,\qquad B_{h+1}\le B_h,\qquad
 B_h\asymp(1+h)^{-1/4},\qquad
 m_h\asymp h^{3/4},\qquad m_{h+1}-m_h\asymp B_h.
\]
The constants are independent of \(h\), and \(b_h(x)=b_h(-x)\).
\end{theorem}
The exact identity and monotonicity, rather than the bridge power alone,
will control convex cuts and variation of the endpoint kernel.

For an internal honeycomb mid-edge \(p\) in a convex domain \(D\),
let \(L_p(D)\) be the sum over ordered boundary chords through \(p\),
weighted by \(\rho^{|\gamma|}\cos(3W(\gamma)/8)\). A polygon nest
around a face center \(f\) is a finite set of pairwise vertex-disjoint
simple polygons, each strictly enclosing \(f\). Its weight is the
product of \(2\rho^{|\ell|}\) over its polygons, and the empty nest
has weight one. Write \(Z_f(D)\) for the nesting mass in \(D\).
Let \(P_2(r)\) denote the whole-plane nesting mass with every polygon's
diameter at most \(r\), and let \(S_k\) denote the middle-face nesting
mass in a strip of \(2k\) bands.
\begin{theorem}[Length and nesting]\label{lem:02-length-nests}\label{thm:05-nest-normalization}
For the two face centers \(f_1,f_2\) adjacent to \(p\),
\[
 \frac1{\sqrt2}\bigl(Z_{f_1}(D)+Z_{f_2}(D)\bigr)
 \le L_p(D)\le Z_{f_1}(D)+Z_{f_2}(D).
\]
The same comparison holds with finite quantities in every infinite strip
of fixed positive height. Moreover,
\(P_2(r)=r^{1/12+o(1)}\) and \(S_k=k^{1/12+o(1)}\).
\end{theorem}
These are the length--nest and planar-normalization theorems of
\cite[Lemma 11.4, Theorem 11.2 and Corollary 11.3]{CylinderCompanion}. We will place diameter-truncated nests in the
interior of a finite triangle; an untruncated planar partition function
would not give that inclusion.
\begin{theorem}[Raw strip first-length mass]\label{prop:02-length-reduction}
From one fixed bottom port of the \(h\)-band strip,
\[
 \sum_{\gamma\text{ arch or bridge}}\rho^{|\gamma|}|\gamma|
       \le h^{13/12+o(1)}.
\]
\end{theorem}
The sum is unnormalized and includes paths ending on either boundary.
It is Proposition 11.5 of \cite{CylinderCompanion}; it is
stronger information for our purpose than a mean conditional on reaching
a deep line. The main geometric proof will supply a matching lower
first moment on bridges confined to diameter at most a fixed multiple
of their height.

\subsection{Marked polygons and arcs}
The marked-polygon theorem of \cite[Theorem 1.1]{MarkedCompanion} gives
\begin{theorem}[Polygon length-square mass]\label{thm:07_magnetic_analysis-moment}
For unrooted, unoriented simple honeycomb polygons modulo triangular-lattice
translations,
\[
 \sum_{[\ell]:\,\operatorname{diam}\ell\le H}
             \rho^{|\ell|}|\ell|^2\le H^{2/3+o(1)}.
\]
\end{theorem}
Here \(|\ell|\) is its vertex count. The translation-class convention
is needed when an exterior closure is translated back to a rooted path.

We also use the cylinder arc estimate of \cite[Theorem 8.3]{MarkedCompanion}.
Its cylinder is
formed by a row cut in cyclic order \([p,X,p',Y]\), where \(p,p'\)
are two marked ports and \(X,Y\) are sets of ordinary sites. Put
\(s=|X|\), \(b=|Y|\), and require \(s+b+2\) even. The physical row
uses rhombus parameters \(\lambda=\pi/8\) or \(2\lambda\): these
two parameter values label the two lattice side directions forming the
positive cone. After a
common shift the logarithmic parameters \(a_i\) have imaginary heights
\(0\) or \(\beta=\pi/4\) in \(X\) and at the marks, and height
zero in \(Y\). At these physical parameters all real parts of the
shifted \(a_i\) are zero. An upper site is one of height \(\beta\). Let
\(A(p,p')\) be the unnormalized critical mass of self-avoiding cylinder
arcs departing west from \(p\) and arriving east at \(p'\), making no
other visit to either mark and leaving their unused half-edges vacant.
\begin{theorem}[Uniform one-arc bound]\label{09_fast_paths:one-arc-bound}
There are constants \(\delta_1>0\), \(L_*>0\), and \(C<\infty\)
such that \(A(p,p')\le b^C\) whenever
\[
 L_Y\ge L_X\ge L_*,\qquad s\le b,\qquad
 L_X=\tfrac34\log s,\quad L_Y=\tfrac34\log b,
\]
and at least fraction \(1/2-\delta_1\) of the ordinary \(X\)-sites
are upper. There is no positive lower bound on \(L_X/L_Y\).
\end{theorem}
The exponent \(C\) is fixed throughout this domain. In the next section
we construct such a row cut from an arbitrary walk in a box. There
\(s\) may stay bounded while \(b\) grows with the box, which explains
why a comparable-wall estimate would not suffice.

\subsection{A tail for every polygon count}
The final input is different from the length-square bound. On the
balanced physical two-marker cylinder with even \(N=2m\), use the
parameter \(2\lambda\) on its first \(m\) bands and \(\lambda\)
on the remaining \(m\) bands. Let \(w_N(a)\) be the critical mass of
families of exactly \(a\) disjoint separating polygons, namely
\[
 w_N(a)=\sum_{\substack{\mathcal F\text{ disjoint, separating}\\
                         \#\mathcal F=a}}
               \prod_{P\in\mathcal F}\rho^{|P|}.
\]
There is no factor two per polygon in this definition. The all-count
tail bound of \cite[Corollary 8.2]{CylinderCompanion} states that, uniformly for
every integer \(a\ge0\),
\begin{equation}\label{eq:input-cylinder-tail}
 w_N(a)\le \exp\left(C\log N-\frac{ca^2}{\log N}\right)
\end{equation}
for sufficiently large even \(N\). Its proof uses the all-real-fugacity estimate of
\cite[Theorem 8.1]{CylinderCompanion}, namely the bound
\(Z_N(2\cosh v)\le\exp(C\log N(1+v^2))\) for every real \(v\ge0\).
An asymptotic for bounded fugacity would not justify arbitrarily large
\(a\). We use \eqref{eq:input-cylinder-tail} only after an explicit
reflection and concatenation has turned prescribed slab crossings into
such polygon families.
